\chapter{功与能}

在牛顿力学中，力描述的是\emph{状态如何改变}，而能量描述的是\emph{改变能够以怎样统一的数量形式被记录}。前几章我们已经用微积分语言建立了位移、速度、加速度与动量的关系；本章进一步说明：当力沿着位移方向持续作用时，它通过\emph{做功}把机械运动的变化压缩为一个标量关系，并自然导出动能、势能、机械能守恒与功率等核心概念。

从高中物理的角度看，“功与能”常常是一组公式；从微积分的角度看，它们是一条清晰的逻辑链：
\[
\vb{F}(\vb{r}) \longrightarrow W=\int_C \vb{F}\cdot \dd \vb{r}
\longrightarrow W=\Delta E_k
\longrightarrow U(\vb{r})=-\int \vb{F}\cdot \dd \vb{r}
\longrightarrow E_k+U=\text{常量}.
\]
本章将围绕这条链逐步展开。

\section{功的定义}\label{sec:work-definition}

\begin{definition}[力做功]
若质点在外力 $\vb{F}$ 作用下沿曲线 $C$ 从点 $A$ 运动到点 $B$，取轨道上一小段位移 $\dd \vb{r}$，则该小段上的元功定义为
\[
\dd W = \vb{F}\cdot \dd \vb{r}.
\]
总功定义为沿轨道的线积分
\[
W_{A\to B}=\int_C \vb{F}\cdot \dd \vb{r}.
\]
\end{definition}

这个定义直接揭示了三点：
\begin{enumerate}
  \item 功是标量，由力和位移的点积给出；
  \item 只有力在位移方向上的分量才做功；
  \item 对变力或曲线运动，必须通过积分把每一小段元功累加起来。
\end{enumerate}

若运动是一维直线运动，位置坐标为 $x$，力为 $F(x)$，则
\[
W=\int_{x_1}^{x_2} F(x)\,\dd x.
\]
于是，在 $F-x$ 图像下，功就是曲线与 $x$ 轴围成的\emph{带符号面积}：图像在轴上方时功为正，在轴下方时功为负。

\begin{remark}
“功”描述的是力转移能量的过程，而不是力本身的大小。大力未必做大功：若位移很小，或力始终与位移垂直，则功可以很小甚至为零。例如匀速圆周运动中的向心力始终垂直于瞬时位移，因此不做功。
\end{remark}

\begin{example}[恒力做功与夹角]
一个大小恒为 $F$ 的力与位移 $\vb{s}$ 的夹角为 $\theta$，求该力做功。

\textbf{解：}
恒力情况下，线积分可直接提出常量：
\[
W=\int_C \vb{F}\cdot \dd \vb{r}=\vb{F}\cdot\int_C \dd \vb{r}=\vb{F}\cdot \vb{s}=Fs\cos\theta.
\]
这就是高中熟悉的公式，但现在我们知道它是线积分在恒力情形下的特例。
\end{example}

\begin{example}[变力做功：线性增大的拉力]\label{ex:linear-force-work}
某物体沿 $x$ 轴由 $x=0$ 移到 $x=a$，受力 $F(x)=kx$，其中 $k>0$ 为常量。求该力做功。

\textbf{解：}
由一维积分公式，
\[
W=\int_0^a kx\,\dd x = \frac{1}{2}ka^2.
\]
这说明：当力随位移线性增大时，平均力为末态力的一半，因此做功等于三角形面积。
\end{example}

\begin{figure}[htbp]
\centering
\begin{tikzpicture}
\begin{axis}[
  width=0.75\textwidth, height=0.5\textwidth,
  xlabel={位置 $x$}, ylabel={力 $F$},
  axis lines=middle,
  grid=major, grid style={gray!30},
  thick,
  xmin=0, xmax=4.5, ymin=0, ymax=3.2,
  xtick={4}, xticklabels={$a$},
]
\addplot[name path=force, blue, domain=0:4, samples=100] {0.7*x};
\addplot[name path=axisline, draw=none, domain=0:4] {0};
\addplot[blue!18] fill between[of=force and axisline];
\end{axis}
\end{tikzpicture}
\caption{变力做功的图像意义：$F(x)$ 曲线下的面积就是功}
\end{figure}

\begin{example}[二维路径上的线积分]
平面内有力场
\[
\vb{F}(x,y)=(-ky,\,kx),
\]
质点从原点沿直线 $y=x$ 运动到点 $(a,a)$。求该力做功。

\textbf{解：}
参数方程取为
\[
\vb{r}(t)=(t,t),\qquad 0\le t\le a,
\]
则
\[
\dd \vb{r}=(1,1)\,\dd t,
\qquad
\vb{F}(\vb{r}(t))=(-kt,kt).
\]
因此
\[
\vb{F}\cdot \dd \vb{r}=(-kt,kt)\cdot (1,1)\,\dd t = 0.
\]
从而
\[
W=\int_0^a 0\,\dd t =0.
\]
这里并不是“没有力”，而是因为该力在整条路径上的切向分量恰好处处相互抵消。
\end{example}

\begin{definition}[变力做功的微积分表述]
若质点运动轨迹写成 $\vb{r}(t)$，其中 $t_1\le t\le t_2$，则线积分可改写为
\[
W=\int_{t_1}^{t_2} \vb{F}(\vb{r}(t))\cdot \dv{\vb{r}}{t}\,\dd t
=\int_{t_1}^{t_2} \vb{F}\cdot \vb{v}\,\dd t.
\]
它表明：做功是“力与速度的点积”对时间的累积。
\end{definition}

\begin{remark}
以后我们定义功率 $P=\dv*{W}{t}$ 时，将立即得到 $P=\vb{F}\cdot \vb{v}$。因此，功与功率并不是两个彼此孤立的公式，而是积分与导数的关系。
\end{remark}

\section{动能定理}\label{sec:work-kinetic}

力做功为什么如此重要？根本原因在于：它直接决定速度大小的变化，这就是动能定理。

\begin{definition}[动能]
质量为 $m$、速率为 $v$ 的质点，其动能定义为
\[
E_k=\frac{1}{2}mv^2.
\]
\end{definition}

\begin{law}[动能定理]
合外力对质点所做的总功，等于质点动能的增量：
\[
W_{\text{合}} = \Delta E_k = E_{k2}-E_{k1}.
\]
\end{law}

\begin{proof}
由牛顿第二定律
\[
\vb{F}_{\text{合}}=m\dv{\vb{v}}{t}.
\]
两边与位移元 $\dd \vb{r}=\vb{v}\,\dd t$ 点乘：
\[
\vb{F}_{\text{合}}\cdot \dd \vb{r}
= m\dv{\vb{v}}{t}\cdot \vb{v}\,\dd t.
\]
左边就是元功 $\dd W$。右边利用恒等式
\[
\dv{(v^2)}{t}=\dv{(\vb{v}\cdot \vb{v})}{t}=2\vb{v}\cdot\dv{\vb{v}}{t},
\]
得
\[
\dd W = m\vb{v}\cdot\dv{\vb{v}}{t}\,\dd t
= \frac{m}{2}\dd(v^2)
=\dd\qty(\frac{1}{2}mv^2).
\]
对初态到末态积分，得到
\[
W_{\text{合}}=\int \dd W
=\int \dd\qty(\frac{1}{2}mv^2)
=\frac{1}{2}mv_2^2-\frac{1}{2}mv_1^2
=\Delta E_k.
\]
证毕。
\end{proof}

这个证明的关键不是技巧，而是观念转换：
\[
\text{牛顿第二定律 } \vb{F}=m\vb{a}
\quad \Longrightarrow \quad
\text{做功关系 } \vb{F}\cdot \dd\vb{r}=\dd E_k.
\]
前者是矢量微分方程，后者是标量积分关系。许多复杂运动问题，用后者往往比直接解加速度更简洁。

\begin{example}[动能定理求末速度]
质量为 $m$ 的物块在水平面上受恒力 $F$ 作用前进距离 $s$，摩擦力大小恒为 $f$，初速度为 $v_0$。求末速度 $v$。

\textbf{解：}
合外力做功为
\[
W=(F-f)s.
\]
由动能定理，
\[
(F-f)s=\frac{1}{2}mv^2-\frac{1}{2}mv_0^2.
\]
故
\[
v^2=v_0^2+\frac{2(F-f)s}{m}.
\]
当右端为零时物块恰好停下；若右端为负，则说明物块在走完 $s$ 之前已经停止，此时题目所设全过程并不成立。
\end{example}

\begin{example}[由动能定理反求变力]\label{ex:reverse-force}
某质点沿 $x$ 轴运动，质量为 $m$，速度满足
\[
v^2=v_0^2+2\alpha x,
\]
其中 $\alpha$ 为常量。求合外力 $F(x)$，并说明其做功规律。

\textbf{解：}
由动能
\[
E_k=\frac{1}{2}mv^2=\frac{1}{2}m(v_0^2+2\alpha x)
=\frac{1}{2}mv_0^2+m\alpha x.
\]
因此动能随位移线性增加，故有
\[
\dv{E_k}{x}=m\alpha.
\]
又因一维情形下
\[
F(x)=\dv{W}{x}=\dv{E_k}{x},
\]
故
\[
F(x)=m\alpha,
\]
即合外力为恒力。并且从 $0$ 到 $x$ 的总功为
\[
W=m\alpha x.
\]
这道题说明：若已知 $v^2$ 与 $x$ 的函数关系，常常可以不经时间，直接借助动能定理研究力学量。
\end{example}

\begin{remark}
动能定理适用于任意力，只要讨论的是\emph{合外力总功}，就一定成立；而机械能守恒只在满足特定条件时才成立。因此从适用范围看，动能定理比机械能守恒更普遍。
\end{remark}

\section{势能}\label{sec:potential-energy}

并不是所有力都能由“位置”唯一决定其做功结果。若某种力对物体从 $A$ 到 $B$ 做的功只与始末位置有关、与路径无关，我们就称它为保守力。

\begin{definition}[保守力]
若力 $\vb{F}$ 在某一区域内满足：质点从点 $A$ 到点 $B$ 的做功与路径无关，只由始末位置决定，则称 $\vb{F}$ 为\emph{保守力}。等价地，对任意闭合路径 $C$ 有
\[
\oint_C \vb{F}\cdot\dd\vb{r}=0.
\]
\end{definition}

\begin{definition}[势能函数]
若 $\vb{F}$ 为保守力，则存在标量函数 $U(\vb{r})$，称为势能，使得
\[
W_{A\to B}^{(\vb{F})}=U(A)-U(B)=-(U(B)-U(A)).
\]
在一维情形中，可写为
\[
U(x)=-\int F(x)\,\dd x + C,
\]
其中 $C$ 是任意常数，反映势能零点可以自由选取。
\end{definition}

由上式立刻得到一维关系
\[
F(x)=-\dv{U}{x},
\]
多维情形可写成
\[
\vb{F}=-\grad U.
\]
这说明：力总是指向势能降低最快的方向。

\begin{theorem}[保守力做功与势能变化]
对任意保守力，沿任一路径从 $A$ 到 $B$ 的做功满足
\[
W_{A\to B}^{(\vb{F})}=-\Delta U.
\]
因此，保守力做正功时势能减少；保守力做负功时势能增加。
\end{theorem}

\subsection*{引力势能}

在近地面重力场中，重力近似恒为 $\vb{G}= -mg\uvec{y}$。若竖直坐标取向上为正，则
\[
F_y=-mg,
\qquad
U(y)=-\int (-mg)\,\dd y = mgy + C.
\]
通常取地面某处为零势能面，则
\[
U_g=mgh.
\]
更一般地，对万有引力
\[
F(r)=-\frac{GMm}{r^2},
\]
势能为
\[
U(r)=-\int F(r)\,\dd r = -\frac{GMm}{r}+C.
\]
若规定无穷远处势能为零，则
\[
U(r)=-\frac{GMm}{r}.
\]

\subsection*{弹性势能}

对理想轻弹簧，胡克定律给出回复力
\[
F(x)=-kx,
\]
其中 $x$ 为相对平衡位置的形变量。于是
\[
U(x)=-\int (-kx)\,\dd x=\frac{1}{2}kx^2+C.
\]
若取平衡位置 $x=0$ 处势能为零，则
\[
U_s=\frac{1}{2}kx^2.
\]

\begin{example}[由力求势能]
一维力场满足
\[
F(x)=-2ax+b,
\]
其中 $a>0$，$b$ 为常量。求势能函数，并指出平衡位置。

\textbf{解：}
由定义
\[
U(x)=-\int (-2ax+b)\,\dd x = ax^2-bx+C.
\]
平衡位置由 $F(x)=0$ 决定，即
\[
-2ax+b=0
\quad\Longrightarrow\quad
x_0=\frac{b}{2a}.
\]
把势能配方，得
\[
U(x)=a\qty(x-\frac{b}{2a})^2 + C-\frac{b^2}{4a}.
\]
因为二次项系数为正，故该平衡位置是稳定平衡点。
\end{example}

\begin{example}[弹簧势能与做功]
轻弹簧劲度系数为 $k$，从原长缓慢压缩到形变量 $x$。求外力所做的功与弹簧势能的变化。

\textbf{解：}
若缓慢压缩，则任意时刻外力大小与弹簧弹力大小相等，即
\[
F_{\text{外}}(s)=ks.
\]
于是外力做功为
\[
W_{\text{外}}=\int_0^x ks\,\dd s=\frac{1}{2}kx^2.
\]
这全部转化为弹簧势能增量：
\[
\Delta U_s = \frac{1}{2}kx^2.
\]
而弹簧本身对外做功则为
\[
W_{\text{弹}}=-\frac{1}{2}kx^2.
\]
正负号恰好相反。
\end{example}

\begin{remark}
势能不是物体“自己带着的一种孤立属性”，而是系统在特定相互作用下由位置决定的能量。例如重力势能属于“地球—物体”系统，弹性势能属于“弹簧—物体”系统。
\end{remark}

\section{机械能守恒}\label{sec:mechanical-energy}

动能定理讨论的是“所有力的总功”，势能讨论的是“保守力的功”。二者结合，就得到机械能守恒定律。

\begin{definition}[机械能]
动能与势能之和称为机械能：
\[
E=E_k+U.
\]
\end{definition}

\begin{law}[机械能守恒定律]
若系统内只有保守力做功，或非保守力总功为零，则系统机械能保持不变：
\[
E_k+U=\text{常量}.
\]
\end{law}

\begin{proof}
由动能定理，系统所受各力总功满足
\[
W_{\text{总}}=\Delta E_k.
\]
把总功分解为保守力功与非保守力功：
\[
W_{\text{保}}+W_{\text{非保}}=\Delta E_k.
\]
而保守力功满足
\[
W_{\text{保}}=-\Delta U.
\]
代入得
\[
-\Delta U + W_{\text{非保}} = \Delta E_k.
\]
移项可得
\[
\Delta(E_k+U)=W_{\text{非保}}.
\]
因此，当 $W_{\text{非保}}=0$ 时，
\[
\Delta(E_k+U)=0,
\]
即
\[
E_k+U=\text{常量}.
\]
证毕。
\end{proof}

\begin{figure}[htbp]
\centering
\begin{tikzpicture}
\begin{axis}[
  width=0.75\textwidth, height=0.5\textwidth,
  xlabel={位置 $x$}, ylabel={能量},
  axis lines=middle,
  grid=major, grid style={gray!30},
  thick,
  xmin=-3.4, xmax=3.4, ymin=0, ymax=3.0,
  legend style={at={(0.98,0.98)}, anchor=north east, draw=none, fill=white},
]
\addplot[blue, domain=-3.2:3.2, samples=100] {0.25*x^2 + 0.4};
\addplot[red, domain=-2.95:2.95, samples=100] {2.6 - (0.25*x^2 + 0.4)};
\addplot[green!50!black, domain=-3.2:3.2, samples=2, dashed] {2.6};
\legend{$U(x)$,$E_k(x)$,$E$}
\end{axis}
\end{tikzpicture}
\caption{机械能守恒时，势能升高的部分正对应动能降低的部分，总能量 $E$ 保持不变}
\end{figure}

这个证明非常重要，因为它把“是否守恒”的判断压缩成一句话：
\[
\boxed{\ \text{机械能是否守恒，取决于非保守力是否做功。}\ }
\]

\begin{theorem}[机械能变化定理]
更一般地，即使机械能不守恒，也总有
\[
\Delta E = \Delta(E_k+U)=W_{\text{非保}}.
\]
它说明非保守力的功负责把机械能转化为其他形式的能，如内能、电能或化学能。
\end{theorem}

\begin{remark}[判断机械能守恒的条件]
判断时可依次检查：
\begin{enumerate}
  \item 是否选定了明确的系统；
  \item 系统内涉及的相互作用中，哪些力可视为保守力；
  \item 摩擦力、阻力、外界拉力、发动机牵引力等非保守力是否做功；
  \item 若有支持力、拉力等非保守力，它们是否恰好与位移垂直，从而做功为零。
\end{enumerate}
例如光滑斜面上的支持力虽然不是保守力，但因始终垂直于位移而不做功，因此不破坏机械能守恒。
\end{remark}

\begin{example}[自由下落的机械能守恒]
质量为 $m$ 的物体由高度 $h_1$ 处从静止释放，下落到高度 $h_2$。忽略空气阻力，求速度。

\textbf{解：}
系统只受重力作用，机械能守恒：
\[
\frac{1}{2}mv_1^2+mgh_1=\frac{1}{2}mv_2^2+mgh_2.
\]
由 $v_1=0$，得
\[
\frac{1}{2}mv_2^2=mg(h_1-h_2),
\qquad
v_2=\sqrt{2g(h_1-h_2)}.
\]
结果与质量无关，这正是自由落体结论的能量版表述。
\end{example}

\begin{example}[含摩擦时机械能不守恒]
物块从粗糙斜面顶端由静止下滑，高度差为 $h$，斜面长为 $\ell$，动摩擦因数为 $\mu$，倾角为 $\theta$。求到达底端时速度。

\textbf{解：}
重力势能减少 $mgh$，摩擦力做负功
\[
W_f=-\mu mg\cos\theta\,\ell.
\]
由机械能变化定理，
\[
\frac{1}{2}mv^2 - 0 + (0-mgh)=W_f,
\]
即
\[
\frac{1}{2}mv^2 = mgh-\mu mg\cos\theta\,\ell.
\]
故
\[
v=\sqrt{2g\qty(h-\mu \ell\cos\theta)}.
\]
这里机械能没有守恒，但“机械能变化等于非保守力做功”仍然成立。
\end{example}

\section{功率}\label{sec:power}

功告诉我们“总共转移了多少能量”，而功率告诉我们“转移得有多快”。

\begin{definition}[平均功率与瞬时功率]
在时间间隔 $\Delta t$ 内，平均功率定义为
\[
\bar{P}=\frac{\Delta W}{\Delta t}.
\]
瞬时功率定义为
\[
P=\dv{W}{t}.
\]
\end{definition}

由元功公式 $\dd W=\vb{F}\cdot\dd\vb{r}$，立刻得到
\[
P=\dv{W}{t}=\vb{F}\cdot\dv{\vb{r}}{t}=\vb{F}\cdot\vb{v}.
\]

\begin{law}[功率公式]
瞬时功率等于力与速度的点积：
\[
P=\vb{F}\cdot \vb{v}=Fv\cos\theta.
\]
\end{law}

当力与速度同向时，功率为正；反向时为负；垂直时功率为零。

\begin{example}[汽车牵引功率]
汽车在平直公路上以速度 $v$ 匀速行驶，阻力大小恒为 $f$。求发动机输出功率。

\textbf{解：}
匀速时牵引力与阻力平衡，故牵引力大小 $F=f$，方向与速度同向。因此
\[
P=Fv=fv.
\]
这说明：同样的阻力下，速度越大，维持匀速所需功率越大。
\end{example}

\begin{remark}
许多“额定功率”问题本质上都在使用
\[
F=\frac{P}{v},
\]
即在功率一定时，速度越大，可提供的牵引力越小。这也是汽车起步爬坡与高速巡航时动力表现不同的原因之一。
\end{remark}

\section{能量图}\label{sec:energy-diagram}

当力是保守力时，研究运动不一定非要回到受力分析；直接观察势能曲线 $U(x)$，往往可以更快把握运动性质。

\begin{definition}[能量图]
在一维保守力问题中，把势能表示为位置的函数 $U(x)$，并在同一图上考虑总机械能水平线 $E$，这种分析方法称为\emph{能量图法}。
\end{definition}

由
\[
E=E_k+U(x)=\frac{1}{2}mv^2+U(x)
\]
可知在任意允许位置上必须有
\[
E-U(x)=\frac{1}{2}mv^2\ge 0.
\]
因此：
\begin{enumerate}
  \item 若 $U(x)>E$，则该位置不可达；
  \item 若 $U(x)=E$，则 $v=0$，这类点称为转折点；
  \item 若 $U(x)<E$，则粒子可以在该区域运动。
\end{enumerate}

\begin{figure}[htbp]
\centering
\begin{tikzpicture}
\begin{axis}[
  width=0.75\textwidth, height=0.5\textwidth,
  xlabel={位置 $x$}, ylabel={势能 $U$},
  axis lines=middle,
  grid=major, grid style={gray!30},
  thick,
  xmin=-3.4, xmax=3.4, ymin=0, ymax=3.2,
]
\addplot[blue, domain=-3.2:3.2, samples=100] {0.28*x^2 + 0.5};
\addplot[red, domain=-3.2:3.2, samples=2, dashed] {2.5};
\addplot[only marks, mark=*, mark size=1.8pt, red] coordinates {(-2.67,2.5) (2.67,2.5)};
\addplot[dashed] coordinates {(-2.67,0) (-2.67,2.5)};
\addplot[dashed] coordinates {(2.67,0) (2.67,2.5)};
\node at (axis cs:-2.67,-0.12) {$x_{\mathrm{L}}$};
\node at (axis cs:2.67,-0.12) {$x_{\mathrm{R}}$};
\end{axis}
\end{tikzpicture}
\caption{势能曲线与总能量线的交点给出转折点：粒子到达这些位置时瞬时速度为零}
\end{figure}

\begin{definition}[平衡点与稳定性]
若某点 $x_0$ 处力为零，即
\[
F(x_0)=0,
\]
则称 $x_0$ 为平衡点。由于 $F(x)=-\dv*{U}{x}$，平衡点等价于
\[
\dv{U}{x}\bigg|_{x=x_0}=0.
\]
若在 $x_0$ 附近，势能有局部最小值，则为稳定平衡；若势能有局部最大值，则为不稳定平衡。
\end{definition}

若 $U(x)$ 二阶可导，则可进一步判别：
\[
U'(x_0)=0,
\qquad
\begin{cases}
U''(x_0)>0, & \text{稳定平衡},\\[4pt]
U''(x_0)<0, & \text{不稳定平衡}.
\end{cases}
\]

\begin{center}
\begin{tikzpicture}[scale=0.95]
  \draw[->] (-4.2,0) -- (4.4,0) node[right] {$x$};
  \draw[->] (0,-0.4) -- (0,4.1) node[above] {$U$};
  \draw[thick,blue,domain=-3.6:3.6,samples=200] plot (\x,{0.12*(\x*\x-4)^2+0.6});
  \draw[thick,red,dashed] (-4,2.4) -- (4,2.4) node[right] {$E_1$};
  \draw[thick,green!50!black,dashed] (-4,3.5) -- (4,3.5) node[right] {$E_2$};
  \fill ( -2,0.6) circle (2pt) node[below] {$x_1$};
  \fill ( 0,2.52) circle (2pt) node[above] {$x_2$};
  \fill ( 2,0.6) circle (2pt) node[below] {$x_3$};
  \node[blue!70!black] at (2.8,1.2) {$U(x)$};
  \node at (-3.0,2.7) {束缚运动};
  \node at (2.8,3.75) {越阱运动};
\end{tikzpicture}
\end{center}

图中两个低谷对应稳定平衡点，中央峰顶对应不稳定平衡点。若总能量为 $E_1$，粒子只能困在左阱或右阱内做往复运动；若总能量提高到 $E_2$，则粒子可以跨越势垒。

\begin{example}[由势能图判断运动]
设某粒子的势能函数为
\[
U(x)=ax^2-bx^3,\qquad a,b>0.
\]
求平衡点并判断稳定性。

\textbf{解：}
由
\[
F(x)=-\dv{U}{x}=-(2ax-3bx^2),
\]
平衡点满足
\[
2ax-3bx^2=0
\quad\Longrightarrow\quad
x=0 \quad\text{或}\quad x=\frac{2a}{3b}.
\]
再看二阶导数：
\[
U''(x)=2a-6bx.
\]
所以
\[
U''(0)=2a>0,
\]
故 $x=0$ 为稳定平衡；而
\[
U''\qty(\frac{2a}{3b})=2a-4a=-2a<0,
\]
故 $x=\dfrac{2a}{3b}$ 为不稳定平衡。
\end{example}

\begin{remark}
能量图法特别适合处理“能否到达某位置”“最大速度在哪里”“转折点在何处”“平衡位置是否稳定”等问题，因为这些问题都可以直接从 $E-U(x)$ 的正负看出，而不必先写加速度方程再积分。
\end{remark}

\label{sec:classic-work-energy}

下面选取三类经典题：变力做功、弹簧问题、传送带问题。我们的目标不是重复常规解法，而是展示微积分语言如何把图像、过程与公式统一起来。

\section*{本章小结}

本章的核心关系可以浓缩为下列四行：
\begin{align*}
\dd W &= \vb{F}\cdot \dd\vb{r},\\
W_{\text{合}} &= \Delta E_k,\\
W_{\text{保}} &= -\Delta U,\\
\Delta(E_k+U) &= W_{\text{非保}}.
\end{align*}
其中，动能定理总成立；势能只对保守力有意义；机械能守恒则是“非保守力不做功”时的特殊结果。把这些公式串起来，很多看似不同的题型——下落、滑块、弹簧、圆周、传送带——便会呈现出统一结构。

